\par\noindent\textbf{(1)} 取度规(+−−−)，s=(pa+pb)²=c²(ma²+mb²)+2pb·pa。若用能量平方变量W²=c²s，则固定靶pb=(mbc,0)给W²=ma²c⁴+mb²c⁴+2mbc²Ea，其中Ea为入射总能量（含静能），不能用仅动能替换。
\par\noindent\textbf{(2)} 阈值时末态在质心系无相对动量，W\_thr=(Σ\_i mi)c²，令M\_f=Σ\_i mi，解出 Ea,thr=[M\_f²−ma²−mb²]c²/(2mb)。入射阈动能Ka,thr=Ea,thr−mac²=[M\_f²−(ma+mb)²]c²/(2mb)，若右端负则反应由其他守恒量决定而非该产生阈值。
\par\noindent\textbf{(3)} 两个等质量m束流头对头且每束总能E，在质心系总三动量抵消，故W=2E；同一总实验室能若用固定靶，一部分进入整体质心平动，W≈√(2mc²E)仅按√E增长。这解释高能对撞机比固定靶更有效地达到重粒子产生阈值。
\par\noindent\textbf{(4)} 光子入射时ma=0、Ea=Eγ，对静止质子mb=mp，若末态总静质量M\_f，阈能 Eγ,thr=(M\_f²−mp²)c²/(2mp)。例如产生额外中性粒子质量mπ且末态p+π⁰时为mπc²[1+mπ/(2mp)]，高于简单静质量mπc²，因为末态还需守恒总动量。
